\documentclass[11pt]{article}
\usepackage[margin=1in]{geometry}
\usepackage{amsmath,amssymb}
\usepackage{booktabs}
\usepackage{hyperref}
\usepackage{siunitx}

\title{Independent Replication Report:\\
\emph{Space-Time Continuous Analysis of Waveform Relaxation\\
for the Heat Equation} (Gander \& Stuart, 1998)}
\author{Independent replication (PDE-100 set, rank 77)}
\date{2026-07-02}

\begin{document}
\maketitle

\begin{abstract}
We independently replicate the four principal claims of Gander \& Stuart,
\emph{SIAM J.\ Sci.\ Comput.} \textbf{19}(6):2014--2031, 1998
(DOI \href{https://doi.org/10.1137/S1064827596305337}{10.1137/S1064827596305337}).
A from-scratch \texttt{numpy} backward-Euler overlapping-Schwarz waveform-relaxation (SWR)
solver reproduces the two-subdomain per-double-iteration contraction rate
$\rho = \alpha(1-\beta)/(\beta(1-\alpha))$ to $\sim 0.1\%$ for three overlaps,
confirms mesh robustness (factor $0.4439$ invariant across an $8\times$ refinement of
$\Delta x=\Delta t$), and matches the eight-subdomain Theorem~3.10 upper bound together
with the predicted early stagnation phase. Three independent free-Argo LLM referees
(\texttt{gpt-5.2}, \texttt{gemini-2.5-pro}, \texttt{gpt-4.1}) unanimously score the
outcome REPLICATED. Verdict: \textbf{REPLICATED}.
\end{abstract}

\section{Paper summary}
Waveform relaxation (WR) is a parallel iterative method for time-dependent problems
in which the space-time domain is split and each piece is solved over the
\emph{whole} time interval, iterating until the interface data converges. Classical
WR discretizes in space first and then algebraically splits the resulting ODE system;
for the heat equation this yields linear (unbounded time) and superlinear (bounded time)
convergence, but the rate degrades under mesh refinement. Gander \& Stuart instead
split in the \emph{physical} domain using \emph{overlapping} domain decomposition
(a Schwarz iteration acting in space-time). Their contributions are:
\begin{itemize}
\item A space-time continuous convergence theory. For two overlapping subdomains
$\Omega_1=[0,\beta L]$ and $\Omega_2=[\alpha L, L]$ with $0<\alpha<\beta<1$ of the 1D
heat equation, the interface error contracts linearly on the infinite time interval
with per-double-iteration factor
\[
\rho \;=\; \frac{\alpha(1-\beta)}{\beta(1-\alpha)}
\qquad (\text{Lemma 2.3, Thm 2.4; semidiscrete Thm 2.8}).
\]
\item This rate is \emph{robust to mesh refinement} provided the physical overlap is
held fixed---the key advantage over algebraic-splitting WR.
\item Generalization to $N$ equal-overlap subdomains with rate bounded above by
$1 - 4r(1-r)\sin^2\!\big(\pi/(2(N+1))\big)$ (Thm 3.10), plus an initial stagnation
phase while boundary information propagates inward.
\item Numerical experiments (\S 4, test problem 4.1) confirming the theory.
\end{itemize}

\paragraph{Test problem (paper eq.~4.1), $L=1$.}
\[
u_t = u_{xx} - \exp\!\big(-(t-1)^2 - (x-1/4)^2\big),\quad
0<x<1,\; 0<t<3,
\]
with $u(0,t)=e^{-2t}$, $u(1,t)=e^{-t}$, $u(x,0)=1$.
Discretization: centered second-order finite differences in space with
$\Delta x = 1/(n+1)$, backward Euler in time, $\Delta x = \Delta t = 0.01$.

\section{Claims table}
\begin{tabular}{cllcc}
\toprule
ID & Claim & Type & Testable? & Tested? \\
\midrule
C1 & 2-subdomain SWR error contracts at rate $\rho=\alpha(1-\beta)/(\beta(1-\alpha))$ & thm+num & yes & yes \\
C2 & The rate is robust to mesh refinement at fixed physical overlap & quantitative & yes & yes \\
C3 & $N$ equal-overlap subdomains: rate $\le 1-4r(1-r)\sin^2(\pi/(2(N+1)))$; initial stagnation & thm+num & yes & yes \\
C4 & Larger overlap $\Rightarrow$ faster convergence; \S 4 numerics match theory & quantitative & yes & yes \\
\bottomrule
\end{tabular}

\section{Method (from scratch)}
All code is in \texttt{work/}; no author code was used (none is distributed).
Language: Python 3, \texttt{numpy} only for the solver (\texttt{matplotlib} for figures).
\begin{enumerate}
\item \textbf{Reference full-domain solver} (\texttt{solve\_full}): backward-Euler heat
solver on $[0,1]$ with Dirichlet BCs $g_1,g_2$ at all times, custom Thomas tridiagonal
solve per step---the ``true'' solution against which the DD algorithm is compared.
\item \textbf{Subdomain solver} (\texttt{solve\_subdomain}): same discretization on a
sub-interval, with time-dependent Dirichlet data at each end supplied from the neighbor's
\emph{previous iterate} interface trace over the entire time strip, implementing the
space-time continuous overlapping-Schwarz WR update (eqs.~2.4--2.7, 2.21--2.24, 3.2--3.3).
\item \textbf{Experiment 1} (\texttt{run\_two\_subdomain}): $\Omega_1=[0,\beta]$,
$\Omega_2=[\alpha,1]$, constant-in-time initial guess equal to the IC value; measured
interface error at grid point $b$ vs the full solution per iteration; fit geometric decay
to obtain the per-double-iteration factor and compared to $\rho$. Overlaps
$(\alpha,\beta) \in \{(0.40,0.60),(0.45,0.55),(0.48,0.52)\}$, $\Delta x=\Delta t = 0.01$.
\item \textbf{Experiment 2} (\texttt{run\_N\_subdomain}): $N=8$ equal subdomains with 35\%
overlap snapped to grid; Jacobi-style parallel sweeps; measured max interface error
(proxy for $\|\xi^k\|_\infty$) vs the Thm 3.10 bound.
\item \textbf{Mesh robustness} (\texttt{mesh\_robust.py}, C2): fixed overlap $(0.4,0.6)$,
$\Delta x$ refined from $0.02$ down to $0.0025$.
\item \textbf{Judging}: three independent free-Argo LLM referees
(\texttt{gpt-5.2}, \texttt{gemini-2.5-pro}, \texttt{gpt-4.1}) scored claims C1--C4.
\end{enumerate}
Reproduce:
\begin{verbatim}
cd work && python3 -m venv .venv && . .venv/bin/activate
pip install numpy scipy matplotlib
python swr_heat.py && python mesh_robust.py && python make_figs.py
\end{verbatim}

\section{Results vs.\ paper}

\subsection*{C1 --- Two-subdomain contraction rate (Fig 4.1)}
\begin{tabular}{lcccc}
\toprule
$(\alpha,\beta)$ & overlap & $\rho$ predicted & measured (per double iter) & rel.\ error \\
\midrule
$(0.40,0.60)$ & $0.20$ & $0.4444$ & $0.4439$ & $0.11\%$ \\
$(0.45,0.55)$ & $0.10$ & $0.6694$ & $0.6690$ & $0.06\%$ \\
$(0.48,0.52)$ & $0.04$ & $0.8521$ & $0.8518$ & $0.04\%$ \\
\bottomrule
\end{tabular}\\[4pt]
The measured per-double-iteration contraction factor matches theoretical $\rho$ to
$\sim 0.1\%$ or better across all three overlaps. Error decays cleanly geometrically
(e.g.\ $(0.4,0.6)$: interface error $0.65 \to 1.3\times 10^{-4}$ over 22 iterations).
\textbf{C1 reproduced.} (Figure \texttt{evidence/fig41\_two\_subdomain.png}.)

\subsection*{C2 --- Mesh robustness (headline claim)}
Fixed overlap $(\alpha,\beta)=(0.4,0.6)$, $\rho_{\text{pred}}=0.4444$:\\[2pt]
\begin{tabular}{ccc}
\toprule
$\Delta x$ & $\Delta t$ & measured (per double iter) \\
\midrule
$0.02$   & $0.02$   & $0.4439$ \\
$0.01$   & $0.01$   & $0.4439$ \\
$0.005$  & $0.005$  & $0.4439$ \\
$0.0025$ & $0.0025$ & $0.4439$ \\
\bottomrule
\end{tabular}\\[4pt]
The contraction factor is invariant to four significant figures across an $8\times$ mesh
refinement, exactly confirming the paper's central claim that the overlapping-Schwarz
WR rate is robust to mesh refinement (in contrast to classical algebraic-splitting WR).
\textbf{C2 reproduced.}

\subsection*{C3 --- $N$-subdomain bound and stagnation (Fig 4.2)}
Eight subdomains, 35\% overlap:
\begin{itemize}
\item Predicted upper bound (per double iter) $=0.9726$; measured decay $=0.9327$.
\item Measured $\le$ bound, as Thm 3.10 requires (it is an upper bound, not equality).
\item Observed $\sim 4$ iterations of stagnation (error $\approx 0.99$) before clean
geometric decay---precisely the paper's remark that ``the error stagnates since information
has to be propagated across domains.''
\end{itemize}
\textbf{C3 reproduced} for the tested $(N,r)$ configuration.
(Figure \texttt{evidence/fig42\_eight\_subdomain.png}.)

\subsection*{C4 --- Overlap vs.\ speed}
From the C1 table, larger overlap $\Rightarrow$ smaller $\rho \Rightarrow$ faster
convergence: overlap $0.20$ gives $\rho=0.444$ (fast), overlap $0.04$ gives $\rho=0.852$
(slow). The monotone relationship, together with the tight agreement with theory
throughout, confirms \S 4's numerical conclusions. \textbf{C4 reproduced.}

\section{LLM-judge panel (free Argo, localhost:44497)}
\begin{tabular}{ll}
\toprule
Judge & Verdict \\
\midrule
\texttt{argo:gpt-5.2}         & REPLICATED (C1,C2 supported; C3,C4 supported/consistent) \\
\texttt{argo:gemini-2.5-pro}  & REPLICATED (all four claims supported) \\
\texttt{argo:gpt-4.1}         & REPLICATED (all four claims supported) \\
\bottomrule
\end{tabular}\\[4pt]
Unanimous REPLICATED. Full referee texts in \texttt{evidence/judges/}.
(\texttt{argo:claude-opus-4.8/4.7} were attempted but hit an Argo chat-endpoint
response-serialization quirk; the gpt/gemini panel was used instead, all free endpoints.)

\section{Genuine critique}
The replication is quantitatively tight on the two headline analytic results (C1,C2),
but a scrupulous reading suggests several limitations that we flag honestly rather than
paper over:

\begin{description}
\item[Scope of C3 is thin.]
Only a single $(N,r)=(8,0.35)$ configuration was tested against an \emph{upper} bound. The
measured rate lying below the bound is necessary but not sufficient to say Theorem~3.10 has
been ``verified'' in a strong sense. A sweep over multiple $N$ and $r$ would tighten this
by exposing whether the ratio of measured-to-bound is stable, monotone in $N$, or shows
regimes where the bound becomes loose. The current evidence supports \emph{consistency}
with C3, not a stress test of it.

\item[Discretization matches the paper exactly, which is a strength \emph{and} a weakness.]
Choosing centred second-order finite differences with backward Euler at
$\Delta x = \Delta t = 0.01$ reproduces \S 4 faithfully, but it means we have not probed
whether the same overlap-in-physical-space robustness survives other reasonable
discretizations (Crank--Nicolson, higher-order FD, spectral, larger CFL-like ratios
$\Delta t / \Delta x^2$). The paper's headline claim is about the continuous method, so
alternative discretizations should preserve C2 asymptotically; showing so empirically
would extend the paper rather than reproduce it, and we did not attempt it.

\item[Contraction factor extracted from a scalar summary.]
For C1 and C2, we fit the per-double-iteration factor from the interface error at a single
grid point $b$ (or $\|\cdot\|_\infty$ over interfaces for C3). This matches the paper's
own presentation, but the space-time continuous theory bounds an $L^\infty$-in-time
maximum-principle quantity; a stronger check would evaluate the contraction of
$\sup_{t\in[0,T]}|\text{error}(b,t)|$ across all iterations, not the peak-value fit we
extracted. The agreement to $\sim 0.1\%$ is convincing evidence that the two agree in
practice on this problem, but the metric is a scalar proxy for a functional.

\item[No stress test against algebraic-splitting WR.]
The paper motivates its result by contrast with algebraic-splitting WR, whose rate
degrades with mesh refinement. We did not implement algebraic-splitting WR and compare;
so our C2 result confirms the \emph{positive} claim (rate is mesh-robust) but does not
demonstrate the \emph{comparative} claim (that this is a real improvement over the
alternative). The comparative claim is well established in the literature independently,
so this is a limitation of scope, not evidence against the paper.

\item[Initial guess and stagnation observation.]
The 4-iteration stagnation phase in the 8-subdomain experiment depends on the
constant-in-time initial guess we used (as specified in the paper). Alternative initial
guesses (e.g.\ warm starts from a coarser solve) could shorten or reshape this phase.
Our observation confirms the paper's qualitative remark under its own conditions but does
not characterise how sensitive the stagnation length is to the initial iterate---an aspect
the paper itself treats only qualitatively.

\item[Free-Argo LLM referees are correlated evaluators.]
All three referees are large closed-source models that may share training data and priors.
Their unanimous REPLICATED verdict adds confidence but should not be interpreted as three
independent expert reviews. The core evidence remains the numerical agreement in Tables C1
and C2, which speaks for itself.
\end{description}

\section*{Verdict}
\textbf{REPLICATED.}

\medskip
\noindent\texttt{WAVE\_RESULT set=PDE-100
paper=Gander-Stuart-1998-waveform-relaxation-heat(DOI:10.1137/S1064827596305337,rank77)
verdict=REPLICATED
dir=\textasciitilde/Dropbox/REPLICATE-PROJECT/PDE-Gander-Stuart-waveform-relaxation-heat-1998
one\_line=From-scratch numpy backward-Euler overlapping-Schwarz WR solver reproduces the
2-subdomain contraction rate rho=alpha(1-beta)/(beta(1-alpha)) to \textasciitilde 0.1\% for three overlaps,
confirms mesh-robustness (factor 0.4439 invariant across 8x refinement), and matches the
8-subdomain Thm 3.10 bound with the predicted stagnation phase; 3/3 free-Argo LLM judges
say REPLICATED.}

\end{document}
